\documentclass[10pt,a4paper]{amsart}
\usepackage[margin=19mm]{geometry}
\usepackage{amsmath,amssymb,mathtools}
\usepackage[scr=rsfs,cal=euler]{mathalfa}
\usepackage{xfrac}
\usepackage[unicode,hidelinks,bookmarks=false]{hyperref}
\newcommand{\ie}{i.e.\ }
\newcommand{\cf}{cf.\ }
\newcommand{\resp}{respectively\ }
\newcommand{\Subsection}{\subsection}
\providecommand{\nobreakdash}{\nobreak\hbox{-}}
\setlength{\emergencystretch}{2em}
\multlinegap0pt
\usepackage{bm}
\usepackage{bbm}
\usepackage{dsfont}
\usepackage{xspace}
\usepackage{cancel}
\usepackage{mathtools}
\usepackage{amsthm}
\makeatletter
\@ifundefined{lemma}{\newtheorem{lemma}{Lemma}}{}
\@ifundefined{proposition}{\newtheorem{proposition}[lemma]{Proposition}}{}
\makeatother
\newcommand{\Bigbrac}[1]{\Bigl( #1 \Bigr)}
\newcommand{\CH}{{\mathcal H}}
\newcommand{\CK}{{\mathcal K}}
\newcommand{\E}{{\mathbb E}}
\newcommand{\N}{{\mathbb N}}
\newcommand{\Z}{{\mathbb Z}}
\newcommand{\brac}[1]{\mathopen{}\left( #1 \mathclose{}\right)}
\newcommand{\inv}{^{-1}}
\newcommand{\uk}{{\underline{k}}}
\newcommand{\uu}{{\underline{u}}}
\title[Quantitative concatenation for polynomial box norms]{Quantitative concatenation for polynomial box norms}
\author{Noah Kravitz, Borys Kuca, James Leng}
\date{Source: arXiv:2407.08636v3 (CC BY 4.0)}
\begin{document}
\maketitle

\noindent\emph{Source and licence.} Quantitative concatenation for polynomial box norms. Version arXiv:2407.08636v3 (CC BY 4.0), distributed under the Creative Commons Attribution 4.0 International licence, as recorded in the arXiv licence metadata for this exact version and shown on the abstract page https://arxiv.org/abs/2407.08636v3. Full rights notice: licenses/arxiv-2407.08636-CC-BY-4.0.txt.\par\smallskip

\noindent\textit{Editorial scope.} Counted unit: the Proposition whose source label is \texttt{P: systems of multilinear equations} in the article (source0.tex lines 2194--2205, source label \texttt{P: systems of multilinear equations}), delivered together with the article's own complete original proof of it (source0.tex lines 2248--2399, one \texttt{proof} environment of 152 lines). No mathematics is written, repaired or re-derived: statement and proof are the article's own text line for line. The proof is detached from the statement in the source: the article states the result at lines 2194--2205 and prints its proof at lines 2248--2399, and the proof environment's own header names the statement, so the association is the article's own. Required context retained uncounted: E:set K (lines 2112--2114); E: H\_l (lines 2116--2119); L: bound on pinned solutions (lines 2122--2131). Every definition, display and label of those lines is retained unchanged. Omitted from this package and not redistributed: 1--2111, 2115--2115, 2120--2121, 2132--2193, 2206--2247, 2400--2660. Nothing inside a retained excerpt is omitted or altered: every retained body is delivered exactly as the article's lines are written, so no omission receipt of any kind accompanies this package. Citations: the retained excerpts carry no citation command at all, so no bibliography entry of the article is transcribed and no reference is redistributed. Reference adaptation: none was needed and none was made; every reference inside the delivered unit resolves inside it, so no unresolved reference or citation is produced. Printed slips kept literal: no printed word, symbol, hypothesis or number of the counted statement or of its proof was altered. Layout and class: the article's own class is \texttt{amsart} with packages including amssymb, amsfonts, theoremref, mathrsfs, hyperref, mdwlist, thm-restate, textcomp, enumerate, inputenc, fontenc, eucal, xcolor, tikz; the wrapper is the lane's pinned \texttt{amsart} editorial wrapper at 10pt on A4, which restates the theorem-like environments the retained text needs and adds the article's own macro definitions that the retained text needs (\texttt{\textbackslash Bigbrac}, \texttt{\textbackslash CH}, \texttt{\textbackslash CK}, \texttt{\textbackslash E}, \texttt{\textbackslash N}, \texttt{\textbackslash Z}, \texttt{\textbackslash brac}, \texttt{\textbackslash inv}, \texttt{\textbackslash uk}, \texttt{\textbackslash uu}), with extra wrapper packages (bm, bbm, dsfont, xspace, cancel, mathtools). The delivered numbering is the unit's own. Assets: no retained span contains a figure environment, a graphics-inclusion command, a PDF-page inclusion, a table or a TikZ picture; no image file is copied into this package, the package declares no asset dependency and no image file is redistributed with it. Licence: version arXiv:2407.08636v3 of this article is distributed under the Creative Commons Attribution 4.0 International licence, as recorded in the arXiv licence metadata for this exact version and shown on the abstract page https://arxiv.org/abs/2407.08636v3. A full rights notice is delivered at licenses/arxiv-2407.08636-CC-BY-4.0.txt. Characters: every retained excerpt is pure ASCII, so the delivered engine, XeLaTeX with its default Latin Modern text font, reports no missing character.
\begin{align}\label{E:set K}
        \CK_{t, \ell, r} := \{(k_{li})_{\substack{(l,i)\in[r]\times[\ell]}}\in[t]^{r\ell}:\; 1\leq k_{l1}<\cdots < k_{l\ell}\leq t\;\; \textrm{for\; all}\;\; l\in[r]\}
    \end{align}

    \begin{multline}\label{E: H_l}
        \CH_{l,\eta} := \{(h_{lk_1\cdots k_l})_{k_1, \ldots, k_l\in[t]}\in[\pm H]^{t^l}:\; |h_{lk_1\cdots k_l}-h_{lk''_1\cdots k''_l}|\geq \eta H,\\ \gcd(h_{lk_1\cdots k_l}-h_{lk''_1\cdots k''_l}, h_{lk'_1\cdots k'_l}-h_{lk''_1\cdots k''_l})\leq \eta\inv\\
    \textrm{for\; distinct}\; (k_1, \ldots, k_l),\; (k'_1, \ldots, k'_l),\; (k''_1, \ldots, k''_l)\in[t]^l\}
    \end{multline}

\begin{lemma}\label{L: bound on pinned solutions}
Let $\eta>0$ and $H, M, l, \ell, r, s, t\in\N$ satisfy $H\leq M$,  $3\leq \ell\leq t$, and $l\leq r$. Let $\CK = \CK_{t, \ell, r}$ be as in \eqref{E:set K} and $\CH = \CH_{l, \eta}$ be as in \eqref{E: H_l}.
Then
    \begin{multline*}
        \E_{\substack{h_{lk_1\cdots k_l}\in[\pm H]:\\ k_1, \ldots, k_l\in[t]}}\; \mathbf{1}_{\CH}((h_{lk_1\cdots k_l})_{k_1, \ldots, k_l})\cdot \prod_{i_1, \ldots, i_{l-1}\in[\ell]}\;\prod_{j\in[s]}\;\prod_{\uk\in\CK}\; \prod_{u_{l+1}, \ldots, u_r\in\{0,1\}}\\ \max_{\substack{n_{0}, n_{1}\in\Z}}\;
        \E_{\substack{|m_{i_l}|\leq H^{u_{l+1}+\cdots + u_r} M:\\ i_l\in[\ell]}}\; \mathbf{1}\Bigbrac{\sum\limits_{i_l\in[\ell]}m_{i_l} = n_0}\cdot
        \mathbf{1}\Bigbrac{\sum\limits_{i_l\in[\ell]}h_{lk_{1i_1}\cdots k_{li_l}}m_{i_l} = n_1}\\
        \ll_{l, \ell, r, s, t}\eta^{-2^{r-l+1}s\ell^{l-1}|\CK|}  M^{-2^{r-l+1}s\ell^{l-1}|\CK|} H^{-(r-l+1)2^{r-l}s\ell^{l-1}|\CK|}.
    \end{multline*}
\end{lemma}

\begin{proposition}[Solutions to systems of multilinear equations]\label{P: systems of multilinear equations}
    Let $\eta>0$, let $\ell, r, s, t\in\N$ with $3\leq\ell\leq t$, and let $H, M\in\N$ with $H\leq M$. Let $\CK = \CK_{t, \ell, r}$ be as in \eqref{E:set K}, and for each $l\in[r]$, let $\CH_l:=\CH_{l,\eta}$ be as in \eqref{E: H_l}. Then 
        \begin{multline*}
        \max_{\substack{n_{j\uk\uu}\in\Z:\; j\in[s],\\ \uk\in\CK,\; \uu\in\{0,1\}^r}}\;\E_{\substack{m_{j\uk i_1\cdots i_r}\in[\pm M]:\\ j\in[s],\; \uk\in\CK,\; i_1, \ldots,  i_r\in[\ell]}}
                \E_{\substack{h_{lk_{1}\cdots k_{l}}\in[\pm H]:\\  k_1, \ldots, k_r\in[t],\; l\in[r]}}\; \Bigbrac{\prod_{l\in[r]}\mathbf{1}_{\CH_l}((h_{lk_{1}\cdots k_{l}})_{k_1, \ldots, k_l \in [t]})}\\
        \prod_{j\in[s]}\; \prod_{\substack{\uk\in\CK}}\prod_{\substack{\uu\in\{0,1\}^r
        }} \mathbf{1}\Bigbrac{\sum\limits_{i_1, \ldots, i_r\in[\ell]} h_{1k_{1i_1}}^{u_1} \cdots h_{rk_{1i_1}\cdots k_{ri_r}}^{u_r}m_{j\uk i_1\cdots i_r} = n_{j\uk\uu}}\\
        \ll_{\ell, r, s, t} \eta^{-O_{\ell, r, s,t}(1)}  M^{-2^{r}s|\CK|}H^{-r2^{r-1}s|\CK|}.
    \end{multline*}
(Here, the superscript $u$'s are exponents, not indices.)

\end{proposition}

\begin{proof}[Proof of Proposition~\ref{P: systems of multilinear equations}] We let all implicit constants depend on $\ell, r, s, t$. The proof goes by induction on $r$.  The $r=1$ case is the content of Lemma \ref{L: bound on pinned solutions}.  The induction step is in general quite notationally complicated, so, to make it more palatable, we first present the case $r=2$.

\smallskip
\textbf{The case $r=2$.}
\smallskip

In this case, our claimed bound takes the form

        \begin{multline*}
        \max_{\substack{n_{j\uk00},\; n_{j\uk01},\\ n_{j\uk10},\; n_{j\uk11}\in\Z:\\ j\in[s],\; \uk\in\CK}}\;\E_{\substack{m_{j\uk i_1 i_2}\in[\pm M]:\\ j\in[s],\; \uk\in\CK,\; i_1, i_2\in[\ell]}}
                \E_{\substack{h_{1k_1}, h_{2k_1k_2}\in[\pm H]:\\ k_1, k_2\in[t]}}\; \mathbf{1}_{\CH_1}((h_{1k_1})_{k_1}) \cdot \mathbf{1}_{\CH_2}((h_{2k_1k_2})_{k_1, k_2})\\
        \prod_{j\in[s]}\; \prod_{\substack{\uk\in\CK}}\prod_{\substack{\uu\in\{0,1\}^2
        }} \mathbf{1}\Bigbrac{\sum\limits_{i_1, i_2\in[\ell]} h_{1k_{1i_1}}^{u_1} h_{2k_{1i_1}k_{2i_2}}^{u_2}m_{j\uk i_1 i_2} = n_{j\uk\uu}}\\
        \ll \eta^{-O(1)}  M^{-4s|\CK|} H^{-4s|\CK|}.
    \end{multline*}
    For each $(j,\uk)\in[s]\times \CK$, we have the system of four equations
\begin{align*}
    \sum\limits_{i_1, i_2\in[\ell]} m_{j\uk i_1 i_2} &= n_{j\uk00}\\
    \sum\limits_{i_1, i_2\in[\ell]} h_{1k_{1i_1}} m_{j\uk i_1 i_2} &= n_{j\uk10}\\
    \sum\limits_{i_1, i_2\in[\ell]} h_{2k_{1i_1}k_{2i_2}} m_{j\uk i_1 i_2} &= n_{j\uk01}\\
    \sum\limits_{i_1, i_2\in[\ell]} h_{1k_{1i_1}} h_{2k_{1i_1}k_{2i_2}} m_{j\uk i_1 i_2} &= n_{j\uk11}.
\end{align*}
This system is quadratic in the $h$ variables, and our goal is to decouple it into two separate systems that are linear in $h$. We accomplish this by introducing two dummy variables
\begin{align*}
    t_{i_1 j \uk 0} = \sum_{i_2\in[\ell]}m_{j\uk i_1 i_2}\quad \textrm{and}\quad t_{i_1 j \uk 1} = \sum_{i_2\in[\ell]}h_{2k_{1i_1}k_{2i_2}} m_{j\uk i_1 i_2}
\end{align*}
and renaming $t_{i_1 i_2 j\uk} = m_{j\uk i_1 i_2}$.  Let us mention the motivation for choosing these variables.  The above system of four equations has three different scales: The first is at scale $\asymp M$, the second and third are at scale $\asymp HM$, and the fourth is at scale $\asymp H^2 M$.  In order to apply the two-scale result from Lemma~\ref{L: bound on pinned solutions}, we want to think of the fourth equation as having scale $\asymp H(HM)$, where the $HM$ here ``matches'' the scale of the third equation.  This is precisely the role played by the second set of dummy variables, which is at scale $\asymp HM$.  The first set of dummy variables, at scale $\asymp M$, pertains to the first and second equations, and the two sets of dummy variables are chosen in such a way that they are related at consecutive scales as in Lemma~\ref{L: bound on pinned solutions}. 

Thus, solving the original system is equivalent to solving
\begin{align*}
    \sum_{i_1\in[\ell]} t_{i_1 j \uk 0} = n_{j\uk00}\quad \textrm{and}\quad
    \sum_{i_1\in[\ell]} h_{1k_{1i_1}} t_{i_1 j \uk 0} = n_{j\uk10}\\
    \sum_{i_1\in[\ell]} t_{i_1 j \uk 1} = n_{j\uk01} \quad \textrm{and}\quad
    \sum_{i_1\in[\ell]} h_{1k_{1i_1}} t_{i_1 j \uk 1} = n_{j\uk11},
\end{align*}
together with
\begin{align*}
    \sum_{i_2\in[\ell]}t_{i_1 i_2 j\uk} = t_{i_1 j \uk 0} \quad \textrm{and}\quad
    \sum_{i_2\in[\ell]}h_{2k_{i_1}k_{i_2}}t_{i_1 i_2 j\uk} = t_{i_1 j \uk 1}\quad \textrm{for\; each}\quad i_1\in[\ell],
\end{align*}
where the variables $t_{i_1 j \uk 0}$ and $t_{i_1 j \uk 1}$ have ranges $[\pm \ell M]$ and $[\pm\ell H M]$, respectively.  Converting the expectation over the $m$'s to a sum, we see that our normalized count is bounded by $O(M^{-s \ell^2|\CK|})$ times 
        \begin{align*}
        &\max_{\substack{n_{j\uk00},\; n_{j\uk01},\\ n_{j\uk10},\; n_{j\uk11}\in\Z:\\ j\in[s],\; \uk\in\CK}}\;
                \E_{\substack{h_{1k_1}\in[\pm H]:\\  k_1\in[t]}}\; \mathbf{1}_{\CH_1}((h_{1k_1})_{k_1}) \sum_{\substack{t_{i_1 j\uk 0}\in[\pm\ell M],\\ t_{i_1 j\uk 1}\in[\pm \ell H M]:\\ j\in[s],\; \uk\in\CK,\; i_1\in[\ell]}}\\
        &\qquad\qquad\qquad\qquad\qquad\qquad\Big(\prod_{j\in[s]}\; \prod_{\substack{\uk\in\CK}}\prod_{\substack{\uu\in\{0,1\}^2}} \mathbf{1}\Bigbrac{\sum\limits_{i_1\in[\ell]} h_{1k_{1i_1}}^{u_1} t_{i_1 j\uk u_2} = n_{j\uk\uu}}\Big)\\
        &\qquad\qquad\qquad\times \Big(\E_{\substack{h_{2k_1k_2}\in[\pm H]:\\  k_1, k_2\in[t]}}\; \mathbf{1}_{\CH_2}((h_{2k_1k_2})_{k_1, k_2})
        \sum_{\substack{t_{i_1 i_2 j\uk}\in[\pm M]:\\ j\in[s],\; \uk\in\CK,\; i_1, i_2\in[\ell]}}\\
        &\qquad\qquad\qquad\qquad\qquad\qquad\prod_{i_1\in[\ell]}\prod_{j\in[s]}\; \prod_{\substack{\uk\in\CK}}\prod_{\substack{u_2\in\{0,1\}}} \mathbf{1}\Bigbrac{\sum\limits_{i_2\in[\ell]} h_{2k_{1i_1}k_{2i_2}}^{u_2} t_{i_1 i_2 j\uk} = t_{i_1 j\uk u_2}}\Big);
    \end{align*}
    we will deal first with the system involving the variables $h_{2k_1k_2}$ and then with the system involving the variables $h_{1k_1}$. Note that
    \begin{align*}
        &\sum_{\substack{t_{i_1 i_2 j\uk}\in[\pm M]:\\ j\in[s],\; \uk\in\CK,\; i_1, i_2\in[\ell]}}\prod_{i_1\in[\ell]}\prod_{j\in[s]}\; \prod_{\substack{\uk\in\CK}}\prod_{\substack{u_2\in\{0,1\}}} \mathbf{1}\Bigbrac{\sum\limits_{i_2\in[\ell]} h_{2k_{1i_1}k_{2i_2}}^{u_2} t_{i_1 i_2 j\uk} = t_{i_1 j\uk u_2}}\\
        &\qquad=\prod_{i_1\in[\ell]}\prod_{j\in[s]}\; \prod_{\substack{\uk\in\CK}}\sum_{\substack{t_{i_1 i_2 j\uk}\in[\pm M]:\\ i_2\in[\ell]}}\; \mathbf{1}\Bigbrac{\sum\limits_{i_2\in[\ell]} t_{i_1 i_2 j\uk} = t_{i_1 j\uk 0}}\cdot \mathbf{1}\Bigbrac{\sum\limits_{i_2\in[\ell]} h_{2k_{1i_1}k_{2i_2}} t_{i_1 i_2 j\uk} = t_{i_1 j\uk 1}}.
    \end{align*}
    Thus, we have
    \begin{align*}
        &\E_{\substack{h_{2k_1k_2}\in[\pm H]:\\  k_1, k_2\in[t]}}\; \mathbf{1}_{\CH_2}((h_{2k_1k_2})_{k_1, k_2})
        \sum_{\substack{t_{i_1 i_2 j\uk}\in[\pm M],\\ j\in[s],\; \uk\in\CK,\; i_1, i_2\in[\ell]}}\\
        &\qquad\qquad\qquad\qquad\prod_{i_1\in[\ell]}\prod_{j\in[s]}\; \prod_{\substack{\uk\in\CK}}\prod_{\substack{u_2\in\{0,1\}}} \mathbf{1}\Bigbrac{\sum\limits_{i_2\in[\ell]} h_{2k_{1i_1}k_{2i_2}}^{u_2} t_{i_1 i_2 j\uk} = t_{i_1 j\uk u_2}}\\
        &\qquad\qquad= \E_{\substack{h_{2k_1k_2}\in[\pm H]:\\  k_1, k_2\in[t]}}\; \mathbf{1}_{\CH_2}((h_{2k_1k_2})_{k_1, k_2})\prod_{i_1\in[\ell]}\prod_{j\in[s]}\; \prod_{\substack{\uk\in\CK}}\\ 
        &\qquad\qquad\qquad\qquad\sum_{\substack{t_{i_1 i_2 j\uk}\in[\pm M]:\\ i_2\in[\ell]}}\; \mathbf{1}\Bigbrac{\sum\limits_{i_2\in[\ell]} t_{i_1 i_2 j\uk} = t_{i_1 j\uk 0}}\cdot \mathbf{1}\Bigbrac{\sum\limits_{i_2\in[\ell]} h_{2k_{1i_1}k_{2i_2}} t_{i_1 i_2 j\uk} = t_{i_1 j\uk 1}}.
    \end{align*}
    Lemma \ref{L: bound on pinned solutions} applied with $l=r=2$ allows us to bound this expression by 
    $$O(\eta^{-O(1)} M^{s\ell^2|\CK|-2s\ell|\CK|} H^{-s\ell|\CK|}).$$
    Similarly, the count of solutions to the system involving the variables $h_{1k_1}$ is
\begin{align*}
        &\max_{\substack{n_{j\uk00},\; n_{j\uk01},\\ n_{j\uk10},\; n_{j\uk11}\in\Z:\\ j\in[s],\; \uk\in\CK}}\;
        \E_{\substack{h_{1k_1}\in[\pm H]:\\  k_1\in[t]}}\; \mathbf{1}_{\CH_1}((h_{1k_1})_{k_1}) \sum_{\substack{t_{i_1 j\uk 0}\in[\pm\ell M],\\ t_{i_1 j\uk 1}\in[\pm\ell H M]:\\ j\in[s],\; \uk\in\CK,\; i_1\in[\ell]}}\\
        &\qquad\qquad\qquad\prod_{j\in[s]}\; \prod_{\substack{\uk\in\CK}}\prod_{\substack{\uu\in\{0,1\}^2}} \mathbf{1}\Bigbrac{\sum\limits_{i_1\in[\ell]} h_{1k_{1i_1}}^{u_1} t_{i_1 j\uk u_2} = n_{j\uk\uu}}\\
        &\qquad\leq  \E_{\substack{h_{1k_1}\in[\pm H]:\\  k_1\in[t]}}\; \mathbf{1}_{\CH_1}((h_{1k_1})_{k_1}) \prod_{j\in[s]}\; \prod_{\substack{\uk\in\CK}}\prod_{\substack{u_2\in\{0,1\}}}\\
        &\qquad\qquad\qquad\max_{n_0, n_1\in\Z}\; \sum_{\substack{t_{i_1 j\uk u_2}\in[\pm\ell H^{u_2} M]:\\  i_1\in[\ell]}}\mathbf{1}\Bigbrac{\sum\limits_{i_1\in[\ell]} t_{i_1 j\uk u_2} = n_0}\cdot \mathbf{1}\Bigbrac{\sum\limits_{i_1\in[\ell]} h_{1k_{1i_1}} t_{i_1 j\uk u_2} = n_1},
\end{align*}
and this can be bounded by
$$O(\eta^{-O(1)} M^{2 s\ell |\CK|-4s|\CK|} H^{s\ell|\CK|-4s|\CK|})$$
using the case $l=1, r=2$ of Lemma \ref{L: bound on pinned solutions}. Note that the big centered equations had an  \emph{inequality} rather than an \emph{equality} because we interchanged the $\max$ with the sums and products.

    Combining these two bounds and the normalising factor $O(M^{-s \ell^2|\CK|})$, we bound the overall normalised count by
    \begin{align*}
        &\ll M^{-s \ell^2|\CK|}\cdot \brac{\eta^{-O(1)} M^{s\ell^2|\CK|-2s\ell|\CK|} H^{-s\ell|\CK|}}\cdot \brac{\eta^{-O(1)} M^{2 s\ell|\CK|-4s|\CK|} H^{s\ell|\CK|-4s|\CK|}}\\
        &= \eta^{-O(1)} M^{-4s|\CK|}H^{-4s|\CK|},
    \end{align*}
    as claimed.

\smallskip
\textbf{The case $r>2$.}
\smallskip

We handle the case $r>2$ similarly. For $l\in[r]$, $i_1, \ldots, i_l\in[\ell]$, $j\in[s]$, $\uk\in\CK$, $\uu\in\{0,1\}^r$, we define the auxiliary variable
\begin{align*}
    t_{i_1\cdots i_l j \uk\uu} =\sum\limits_{i_{l+1}, \ldots, i_r\in[\ell]}h_{(l+1)k_{1i_1} \cdots k_{(l+1)i_{l+1}}}^{u_{l+1}}\cdots h_{r k_{1i_1}\cdots k_{ri_r}}^{u_r}m_{j\uk i_1\cdots i_r}
\end{align*}
at scale $H^{u_{l+1}+\cdots+u_r}M$.
Notice that $t_{i_1\cdots i_l j \uk\uu}$ depends on only the coordinates $u_{l+1}, \ldots, u_r$ of $\uu$ and that the $t$'s satisfy the relation
\begin{align}\label{E: recurrence relation on t}
    t_{i_1\cdots i_l j \uk\uu} = \sum_{i_{l+1}\in[\ell]}h_{(l+1)k_{1i_1} \cdots k_{(l+1)i_{l+1}}}^{u_{l+1}} t_{i_1\cdots i_{l+1} j \uk\uu}.
\end{align}
Iterating \eqref{E: recurrence relation on t} through intermediate scales, we see that {for each $j \in [s]$ and $\uk \in \CK$,} we can rewrite
\begin{align*}
        &\sum_{\substack{m_{j\uk i_1\cdots i_r}\in[\pm M]:\\ 
        %j\in[s],\; \uk\in\CK,\; 
        i_1, \ldots,  i_r\in[\ell]}} \mathbf{1}\Bigbrac{\sum\limits_{i_1, \ldots, i_r\in[\ell]} h_{1k_{1i_1}}^{u_1} \cdots h_{rk_{1i_1}\cdots k_{ri_r}}^{u_r}m_{j\uk i_1\cdots i_r} = n_{j\uk\uu}} \\
        &\qquad\qquad\qquad=\sum_{\substack{|t_{i_1 j\uk\uu}|\ll H^{u_2+\cdots + u_r}M:\\ i_1\in[\ell]}} \mathbf{1}\Bigbrac{\sum\limits_{i_1\in[\ell]} h_{1k_{1i_1}}^{u_1} t_{i_1 j\uk\uu}= n_{j\uk\uu}}\\ 
    &\qquad\qquad\qquad\times \sum_{\substack{|t_{i_1i_2 j\uk\uu}|\ll H^{u_3+\cdots +u_r}M:\\ i_1, i_2\in[\ell]}}\;\prod_{i_1\in[\ell]} \mathbf{1}\Bigbrac{\sum\limits_{i_2\in[\ell]} h_{2k_{1i_1}k_{2i_2}}^{u_2}t_{i_1 i_2 j\uk\uu} = t_{i_1 j\uk\uu}}\\
    &\qquad\qquad\qquad\times \cdots \times  \sum_{\substack{|t_{i_1\cdots i_r j\uk\uu}|\leq M:\\ i_1, \cdots, i_r\in[\ell]}}\;\prod_{i_1,\ldots, i_{r-1}\in[\ell]} \mathbf{1}\Bigbrac{\sum\limits_{i_r\in[\ell]} h_{rk_{1i_1}\cdots k_{r i_r}}^{u_r}t_{i_1 \cdots i_r j\uk\uu} = t_{i_1 \ldots i_{r-1} j\uk\uu}},
\end{align*}
where we note that $t_{i_1 \cdots i_r j\uk\uu} = m_{j\uk i_1\cdots i_r}$. Hence our normalized count is bounded by $O(M^{-s\ell^r|\CK|})$ times
\begin{align*}
            &\max_{\substack{n_{j\uk\uu}\in\Z:\; j\in[s],\\ \uk\in\CK,\; \uu\in\{0,1\}^r}}\; \E_{\substack{h_{1k_1}\in[\pm H]:\\ k_1\in[t]}}\; \mathbf{1}_{\CH_1}((h_{1 k_1})_{k_1})\sum_{\substack{|t_{i_1 j\uk\uu}|\ll H^{u_2+\cdots +u_r}M:\; j\in[s],\\ \uk\in\CK,\; u_2, \ldots, u_r\in\{0,1\},\; i_1\in[\ell]}}\\
        &\qquad\qquad\qquad\qquad\Big(\prod_{j\in[s]}\; \prod_{\substack{\uk\in\CK}}\; \prod_{\substack{u_1, \ldots, u_r\in\{0,1\}}}\mathbf{1}\Bigbrac{\sum\limits_{i_1\in[\ell]} h_{1k_{1i_1}}^{u_1} t_{i_1 j\uk\uu}= n_{j\uk\uu}}\Big)
        \\
        &\qquad\qquad\times \E_{\substack{h_{2k_1k_2}\in[\pm H]:\\ k_1,k_2\in[t]}}\; \mathbf{1}_{\CH_2}((h_{2 k_1 k_2})_{k_1, k_2})\; \sum_{\substack{|t_{i_1i_2 j\uk\uu}|\ll H^{u_3+\cdots+u_r}M:\; j\in[s],\\ \uk\in\CK,\; u_3, \ldots, u_r\in\{0,1\},\; i_1, i_2\in[\ell]}}\\
        &\qquad\qquad\qquad\qquad
        \Big(\prod_{i_1\in[\ell]}\; \prod_{j\in[s]}\; \prod_{\substack{\uk\in\CK}}\; \prod_{\substack{u_2, \ldots, u_r\in\{0,1\}}} \mathbf{1}\Bigbrac{\sum\limits_{i_2\in[\ell]} h_{2k_{1i_1}k_{2i_2}}^{u_2}t_{i_1 i_2 j\uk\uu} = t_{i_1 j\uk\uu}}\Big)\\
        &\qquad\qquad\times \cdots \times \E_{\substack{h_{rk_1\cdots k_r}\in[\pm H]:\\ k_1,\ldots, k_r\in[t]}}\; \mathbf{1}_{\CH_r}((h_{r k_1 \cdots k_r})_{k_1, \ldots, k_r})\; \sum_{\substack{|t_{i_1\cdots i_r j\uk\uu}|\leq M:\; j\in[s],\\ \uk\in\CK,\; i_1,\ldots,i_r\in[\ell]}}   \\
        &\qquad\qquad\qquad\qquad
        %i_1, \ldots, i_r\in[\ell],\; j\in[s],\\ \uk\in\CK,\; u_r\in\{0,1\}}} 
        \prod_{i_1,\ldots, i_{r-1}\in[\ell]}\; \prod_{j\in[s]}\; \prod_{\substack{\uk\in\CK}}\; \prod_{\substack{u_r\in\{0,1\}}}
         \mathbf{1}\Bigbrac{\sum\limits_{i_r\in[\ell]} h_{rk_{1i_1}\cdots k_{r i_r}}^{u_r}t_{i_1 \cdots i_r j\uk\uu} = t_{i_1 \cdots i_{r-1} j\uk\uu}}.
\end{align*}
    As in the case $r=2$, we have decoupled our original system of degree-$r$ multilinear equations into $r$ smaller systems, each consisting of a pair of multilinear equations at consecutive scales.  We will next use Lemma \ref{L: bound on pinned solutions} to bound the counts of solutions for these subsystems individually, starting from the system involving the variables $h_{rk_1\cdots k_r}$. The count of solutions to this system can be rephrased as
    \begin{multline*}
\E_{\substack{h_{rk_1\cdots k_r}\in[\pm H]:\\ k_1,\ldots, k_r\in[t]}}\; \mathbf{1}_{\CH_r}((h_{r k_1 \cdots k_r})_{k_1, \ldots, k_r})  \prod_{i_1,\ldots, i_{r-1}\in[\ell]}\; \prod_{j\in[s]}\; \prod_{\substack{\uk\in\CK}}\\
        \sum_{\substack{|t_{i_1\cdots i_r j\uk\uu}|\leq M:\\ i_r\in[\ell]}}\ 
        \prod_{\substack{u_r\in\{0,1\}}}
         \mathbf{1}\Bigbrac{\sum\limits_{i_r\in[\ell]} h_{rk_{1i_1}\cdots k_{r i_r}}^{u_r}t_{i_1 \cdots i_r j\uk\uu} = t_{i_1 \cdots i_{r-1} j\uk\uu}},
    \end{multline*}
    and bounded by $\ll \eta^{-O(1)}M^{s\ell^{r}|\CK|-2s\ell^{r-1}|\CK|}H^{-s\ell^{r-1}|\CK|}$ using the case $l=r$ of Lemma \ref{L: bound on pinned solutions}. Similarly, the count corresponding to the system involving the variables $h_{(r-1)k_1\cdots k_{r-1}}$ can be rephrased as
    \begin{multline*}
\E_{\substack{h_{(r-1)k_1\cdots k_{r-1}}\in[\pm H]:\\ k_1,\ldots, k_{r-1}\in[t]}}\; \mathbf{1}_{\CH_{r-1}}((h_{(r-1) k_1 \cdots k_{r-1}})_{k_1, \ldots, k_{r-1}})  \prod_{i_1,\ldots, i_{r-2}\in[\ell]}\; \prod_{j\in[s]}\; \prod_{\substack{\uk\in\CK}}\; \prod_{\substack{u_{r}\in\{0,1\}}}\\
        \sum_{\substack{|t_{i_1\cdots i_{r-1} j\uk\uu}|\ll H^{u_r} M:\\ i_{r-1}\in[\ell]}}\ 
        \prod_{\substack{u_{r-1}\in\{0,1\}}}
         \mathbf{1}\Bigbrac{\sum\limits_{i_{r-1}\in[\ell]} h_{(r-1)k_{1i_1}\cdots k_{(r-1) i_{r-1}}}^{u_{r-1}}t_{i_1 \cdots i_{r-1} j\uk\uu} = t_{i_1 \cdots i_{r-2} j\uk\uu}}
    \end{multline*}  
    and bounded by
    \begin{align*}
        \ll \eta^{-O(1)}M^{2s\ell^{r-1}|\CK|-s\ell^{r}|\CK|-4s\ell^{r-1}|\CK|}H^{s\ell^{r-1}|\CK|-4s\ell^{r-1}|\CK|}
    \end{align*}
    with the help of the case $l=r-1$ of Lemma \ref{L: bound on pinned solutions}. Continuing in this manner and combining the results, we can bound our initial normalized count by
    \begin{align*}
        &\ll M^{-s\ell^r|\CK|}\cdot \prod_{l\in[r]} \eta^{-O(1)}  M^{2^{r-l}s\ell^l|\CK|-2^{r-l+1}s\ell^{l-1}|\CK|}\\ 
        &\qquad\qquad\qquad\qquad\qquad\qquad H^{(r-l)2^{r-l-1}s\ell^l|\CK| -(r-l+1)2^{r-l} s\ell^{l-1}|\CK|},
    \end{align*}
    where in each $l$-multiplicand, the negative exponents of $\eta, M, H$ come from an application of Lemma \ref{L: bound on pinned solutions} to bound the count of solutions to the system involving the variables $h_{lk_1\cdots k_l}$ and the positive exponents come from counting the admissible values of the variables $t_{i_1\cdots i_l j\uk\uu}$. The telescoping nature of the exponents leads to cancellations: The exponents of $M$ and $H$ cancel out as
    \begin{gather*}
        -s\ell^r|\CK| + (s\ell^r|\CK| - 2s\ell^{r-1}|\CK|) + \cdots + (2^{r-1}s\ell|\CK| - 2^r s|\CK|) = -2^r s|\CK|,\\
        -s\ell^{r-1}|\CK| + (s\ell^{r-1}|\CK| - 4s\ell^{r-2}|\CK|) + \cdots + ((r-1)2^{r-2}s\ell|\CK| - r2^{r-1}s|\CK|) = -r2^{r-1}s|\CK|,
    \end{gather*}
    respectively, yielding the desired final bound
    \begin{align*}
        \ll\eta^{-O(1)} M^{-2^r s|\CK|} H^{-r2^{r-1}s|\CK|}.
    \end{align*} 
    \end{proof}

\end{document}
